\documentclass[11pt,a4paper]{article}

% Importeer configuratie-instellingen
\usepackage{../../LaTeX/styles/exam_config}
\usepackage{../../LaTeX/styles/statistics-macros}
\usetikzlibrary{arrows.meta, backgrounds, bending, calc, decorations.pathreplacing, intersections, positioning, shapes}
\usepgfplotslibrary{fillbetween, statistics}

% % DEFINED COLORS
\definecolor{octahedroncolor}{RGB}{104, 129, 174}
\colorlet{highlightcolor}{green!80!black}
\colorlet{conceptcolor}{cyan!80!black}

% GENERAL COMMANDS
\newcommand{\set}[1]{\left\{#1\right\}}

\newcommand{\comboptproblem}[2]{
    \begin{center}
        \framebox[0.9\width][c]{
            \begin{quote}
                \renewcommand{\arraystretch}{1.5}%
                \begin{tabular}{@{}lp{.7\textwidth}}
                    \textbf{Gegeven:} & #1 \\
                    \textbf{Vind:} & #2
                \end{tabular}
            \end{quote}
        }
    \end{center}
}

\newcommand{\concept}[1]{\textbf{\color{conceptcolor} #1}}

% THEOREMS
\newtheorem{stelling}{Stelling}

% MATH OPERATORS
\DeclareMathOperator{\suchthat}{zodanig\;dat}



\setlength{\arraycolsep}{2pt}
\renewcommand<>\cellcolor[1]{\only#2{\beameroriginal\cellcolor{#1}}}

% Solution Toggle (to show/hide solutions)
\newif\ifshowanswers
% \showanswerstrue % Toggle this to hide / show solutions
\showanswersfalse % Uncomment to hide solutions

% \setbeamertemplate{footline}{}
% \setbeamertemplate{headline}{}

% Aangepaste TikZ-instellingen
\colorlet{currentcolor}{blue}
\colorlet{optimalcolor}{green!50!black}
\colorlet{feasiblecolor}{cyan!60}

\tikzset{
  activity/.style={draw, thick, fill=gray!10, minimum size=6mm, circle},
  arrow/.style={-stealth, thick},
  axis/.style={semithick, -stealth, line cap=round},
  connector/.style={-stealth, thick, draw=black!70},
  constraintdirection/.style={-stealth, very thick, shorten >=-0.3cm},
  critical/.style={red, very thick},
  cross/.style={cross out, draw, minimum size=2*(#1 - \pgflinewidth), inner sep=0pt, outer sep=0pt},
  current/.style={intersect, currentcolor, inner sep=3pt},
  direction/.style={-stealth, orange, ultra thick, dashed},
  feasible/.style={draw=black, ultra thick, fill=feasiblecolor},
  halfspace/.style={fill=blue!20, opacity=0.2},
  hoogtelijn/.style={objective, -, thin, dashed},
  intersect/.style={circle, fill=black, inner sep=1.5pt},
  objective/.style={ultra thick,-stealth, color=red},
  optimum/.style={current, optimalcolor},
  optimumtext/.style={optimalcolor, thick, fill=green!20, align=center},
  tick/.style={thick}
}

\pgfplotsset{
    compat=newest,
    every axis/.style={
        axis x line=middle,
        axis y line=middle,
        enlarge x limits=0.1,
        enlarge y limits=0.1,
        grid=both,
        grid style={line width=.1pt, draw=gray!30},
        tick style={black},
        label style={font=\small},
        tick label style={font=\scriptsize}
    }
}

% TikZ macros
\newcommand{\drawaxes}[2]{% xmin xmax ymin ymax
  \draw[axis] (-0.5,0) -- (#1,0) node[anchor=north west] {$x_1$};
  \draw[axis] (0,-0.5) -- (0,#2) node[anchor=south east] {$x_2$};

  % Axis ticks and labels
  \foreach \x in {1, ..., #1} 
    \draw (\x, -0.1) -- (\x, 0.1) node[below] {\x};
  \foreach \y in {1, ..., #2}
    \draw (-0.1, \y) -- (0.1, \y) node[left] {\y};
  \node (origin) at (0,0) [anchor=north east] {$O$};
}
\graphicspath{{../../LaTeX/figures/}}

\pgfplotsset{compat=1.18}
\usetikzlibrary{shapes.geometric, arrows.meta, positioning, calc, decorations.markings, patterns}

\setlength{\parindent}{0pt}
\setlength{\parskip}{\baselineskip}

\definecolor{QuestionBlue}{HTML}{2563A6}

\newcommand{\courseLogo}{course_logo.png}
\newcommand{\instituteLogo}{inst_logo.png}
\newcommand{\notestitle}{Introduction to Probability}
\newcommand{\courseid}{Course Code}
\newcommand{\instructorname}{Instructor Name}
\newcommand{\classdate}{\today}
\newcommand{\documentlabel}{\textsc Learning task}

\newcommand{\makeclassnotesheader}{%
    \begin{tcolorbox}[
        enhanced,
        colback=QuestionBlue!3,
        colframe=QuestionBlue,
        boxrule=1.1pt,
        arc=3pt,
        left=12pt,
        right=12pt,
        top=12pt,
        bottom=10pt
    ]
    \noindent
    \begin{minipage}[c]{0.16\linewidth}
        \centering
        \IfFileExists{\courseLogo}
        {\includegraphics[width=0.94\linewidth, keepaspectratio]{\courselogo}}
        {\fbox{\scriptsize Course Logo}}
    \end{minipage}%
    \hfill
    \begin{minipage}[c]{0.16\linewidth}
        \centering
        \IfFileExists{\instituteLogo}
        {\includegraphics[width=0.94\linewidth, keepaspectratio]{\institutelogo}}
        {\fbox{\scriptsize Institute Logo}}
    \end{minipage}
    \vspace{5pt}
    \hrule
    \vspace{8pt}
    \begin{tabular*}{\linewidth}{@{\extracolsep{\fill}}lll@{}}
        \textbf{Course ID:} \courseid &
        \textbf{Instructor:} \instructorname &
        \textbf{Date:} \classdate        
    \end{tabular*}
    \end{tcolorbox}
}

\begin{document}
\pagestyle{empty} % Remove this line if page numbers are desired.

\makeclassnotesheader

\section{Learning task 1: the role of statistics in engineering}

We start with two examples from the bachelor Military Systems and Technology.

\subsection*{Maintenance: the reliability of a circuit}
A circuit operates only if there is a path of functional devices from left to right. 
Devices $A$ and $B$ are connected in parallel, device $C$ is in series.

\begin{figure}[h]
    \centering
    \begin{tikzpicture}[scale=0.9, >=stealth]
        % Main line input/output
        \draw[thick] (-0.5, 0) -- (1, 0);
        \draw[thick] (4, 0) -- (5.5, 0);
        \draw[thick] (7.5, 0) -- (9, 0);

        % Parallel split
        \draw[thick] (1, -1.2) -- (1, 1.2);
        \draw[thick] (1, 1.2) -- (2, 1.2);
        \draw[thick] (1, -1.2) -- (2, -1.2);

        % Blocks A and B
        \draw[thick] (2, 0.6) rectangle (3.5, 1.8) node[pos=.5] {$A$};
        \draw[thick] (2, -1.8) rectangle (3.5, -0.6) node[pos=.5] {$B$};

        % Parallel rejoin
        \draw[thick] (3.5, 1.2) -- (4, 1.2);
        \draw[thick] (3.5, -1.2) -- (4, -1.2);
        \draw[thick] (4, -1.2) -- (4, 1.2);

        % Block C
        \draw[thick] (5.5, -0.6) rectangle (7.5, 0.6) node[pos=.5] {$C$};
    \end{tikzpicture}
    \caption*{Figure 1.1 \\ Circuit with three devices}
\end{figure}

For example, the circuit operates if the devices $A$ and $C$ function, but the circuit fails, if device $C$ does not function.
Assume that all devices function on average for $10$ hours. 
What is the probability that the circuit will operate for at least $9$ hours?

We look at this problem in two different ways, first deterministic (i.e. there is no variability), then stochastic (i.e. there is variability).
\begin{itemize}
    \item In the deterministic setting, all devices will function for precisely $10$ hours. The requested probability is therefore equal to $1$.
    \item In the stochastic case, both a simulation using Excel and a calculation (later in this course) result in a probability of approximately $0.26$.
\end{itemize}
See the Excel-file "PMTPS\_STAT\_1\_stochastic.xlsx" or the print screen in Figure 1.2.

In this file, the above circuit is simulated $10.000$ times. 
In each row, three life times for the devices $A$, $B$ and $C$ are generated. 
The life time of the circuit is calculated in column "D".
Next, in column "E", it is determined whether the circuit still functions after $9$ hours. 
In that case, the number $1$ is shown in column "E".

The requested probability is then calculated in cell "E10003".
The magnitude of this probability is partly dependent on the number of simulations. 
Use the button "F9" to examine the sensitivity at this point.

\begin{figure}[h]
    \centering
    \begin{tabular}{|c|c|c|c|c|c|}
        \toprule
            & \textbf{A} & \textbf{B} & \textbf{C} & \textbf{D} & \textbf{E} \\ 
        \midrule
            \textbf{1} & Device A & Device B & Device C & Circuit & $>$ 9 \\ \midrule
            \textbf{9988} & 6,841 & 14,873 & 8,209 & 8,209 & 0 \\ \midrule
            \textbf{9989} & 4,240 & 8,650 & 15,175 & 8,650 & 0 \\ \midrule
            \textbf{9990} & 18,044 & 6,365 & 38,085 & 18,044 & 1 \\ \midrule
            \textbf{9991} & 1,194 & 1,557 & 5,246 & 1,557 & 0 \\ \midrule
            \textbf{9992} & 38,282 & 8,686 & 0,611 & 0,611 & 0 \\ \midrule
            \textbf{9993} & 7,397 & 14,077 & 18,563 & 14,077 & 1 \\ \midrule
            \textbf{9994} & 4,863 & 1,553 & 14,752 & 4,863 & 0 \\ \midrule
            \textbf{9995} & 6,554 & 5,390 & 2,647 & 2,647 & 0 \\ \midrule
            \textbf{9996} & 6,162 & 32,267 & 10,679 & 10,679 & 1 \\ \midrule
            \textbf{9997} & 3,081 & 1,041 & 11,934 & 3,081 & 0 \\ \midrule
            \textbf{9998} & 15,068 & 5,047 & 97,078 & 15,068 & 1 \\ \midrule
            \textbf{9999} & 3,645 & 0,658 & 5,429 & 3,645 & 0 \\ \midrule
            \textbf{10000} & 9,057 & 14,896 & 0,006 & 0,006 & 0 \\ \midrule
            \textbf{10001} & 3,125 & 17,618 & 7,873 & 7,873 & 0 \\ \midrule
            \textbf{10002} & Mean A & Mean B & Mean C & Mean Circuit & Probability \\ \midrule
            \textbf{10003} & 9,983 & 10,122 & 9,876 & 6,644 & 0,2638 \\ 
        \bottomrule
        \end{tabular}
    \caption*{Figure 1.2 \\ Life length of a circuit}
\end{figure}

\subsection*{Operations Research: Search and Detection}

\begin{figure}[h]
\centering
\begin{tikzpicture}[scale=0.8, >=stealth]
    % Boundaries
    \draw[thick] (0,0) -- (0,5);
    \draw[thick] (6,0) -- (6,5);
    
    % Hatching on edges
    \foreach \y in {0.5, 1, 1.5, 2, 2.5, 3, 3.5, 4, 4.5} {
        \draw (0,\y) -- (-0.3,\y+0.2);
        \draw (6,\y) -- (6.3,\y+0.2);
    }
    
    % Edge Labels
    \node[left] at (-0.5,2.5) {left edge};
    \node[right] at (6.5,2.5) {right edge};
    
    % Target Arrows
    \foreach \x in {1, 1.8, 2.6, 3.4, 4.2, 5} {
        \draw[->, thick] (\x, 4.8) -- (\x, 3.8);
    }
    \draw[->, thick] (2.6, 4.8) -- (2.6, 3.2);
    \draw[->, thick] (4.2, 4.8) -- (4.2, 2.8);
    
    \node at (2, 3.5) {targets};
    \node at (3.5, 3.5) {$U$};

    % Observer horizontal arrows
    \draw[very thick, ->] (0, 1.5) -- (2.3, 1.5);
    \draw[very thick, ->] (5, 1.5) -- (3.7, 1.5);
    \draw[very thick] (2.3, 1.5) -- (3.7, 1.5);
    
    \node at (1.5, 1) {observer};
    \node at (4.0, 1) {$V$};
    
    % Bottom Curly Brace D
    \draw[thick] (0,0.5) .. controls (0,0.1) and (1.5,0.1) .. (3,0.1) node[below=4pt] {$D$} .. controls (4.5,0.1) and (6,0.1) .. (6,0.5);
\end{tikzpicture}
\caption*{Figure 1.3 \\ Linear barrier (reversing course at the edge)}
\end{figure}

An observer is protecting a lane with width $D$ while moving at speed $V$ through the lane according some fixed pattern. Any target that closes the observer to within his sweep radius $R$ is detected. So the observer's detecting device is binary: the target is either detected or not detected.

We assume that targets intend to traverse this lane southwards. We also assume that the target speed $U$ is constant along its path. A further assumption is that we know the intent and the capabilities of the target. More precisely, we assume that we know its speed. Its position,
however, is unknown, so arbitrary. In stochastic terms: all positions have an equal probability of occurrence.

Since the target knows that a preference for the edges (i.e. positions nearby the edges of the lane are more likely than positions in the centre of the lane) leads to a lower probability of detection, the target chooses its position based on the probability function in Figure 1.4.

% \begin{figure}[h]
%     \centering
%     \begin{tikzpicture}
%         \begin{axis}[
%             width=0.85\textwidth,
%             height=5.5cm,
%             xmin=-12, xmax=12,
%             ymin=0, ymax=0.12,
%             xtick={-12,-10,-8,-6,-4,-2,0,2,4,6,8,10,12},
%             ytick={0.00, 0.05, 0.10},
%             yticklabels={0.00, 0.05, 0.10},
%             grid=both,
%             grid style={line width=.1pt, draw=gray!40},
%             major grid style={line width=.2pt, draw=gray!60},
%             axis lines=middle,
%             xlabel={$x$},
%             ylabel={$f(x)$},
%             every axis x label={at={(axis description cs:0.5,-0.1)},anchor=north},
%             every axis y label={at={(axis description cs:-0.08,0.5)},rotate=90,anchor=south}
%         ]
%         \addplot[red, ultra thick] coordinates {
%             (-10, 0.10)
%             (0, 0.00)
%             (10, 0.10)
%         };
%         \end{axis}
%     \end{tikzpicture}
%     \caption{Probability function}
% \end{figure}

So if $D=20\text{ nm}$, the probability that the target chooses a position between $-5\text{ nm}$ and $+5\text{ nm}$ is equal to $0.25$ (i.e. the area under the graph of the probability function between $x=-5$ and $x=+5$), while the probability that the target chooses a position between $-10\text{ nm}$ and $-5\text{ nm}$ or between $5\text{ nm}$ and $10\text{ nm}$ is equal to 0.75.

What is the probability that the target will be detected?

Using probability, calculus and geometry, it is not difficult to derive the following formula
\[
P_{\text{det}} = \left\{ \frac{4R^3}{3D^3} \cdot \left(1 + \left(\frac{V}{U}\right)^2\right) - \frac{2R^2}{D^2} \cdot \sqrt{1 + \left(\frac{V}{U}\right)^2} + \frac{2R}{D} \right\} \cdot \sqrt{1 + \left(\frac{V}{U}\right)^2}.
\]

\paragraph{Example 1}
\begin{align*}
D &= 20 \\
R &= 4 \\
V &= 18 \\
U &= 12 \\
P_{\text{det}} &= 0.5236
\end{align*}

\end{document}